\documentclass[11pt,a4paper]{article}

\usepackage[utf8]{inputenc}
\usepackage[T1]{fontenc}
\usepackage{amsmath,amssymb,amsthm,mathtools}
\usepackage{geometry}
\usepackage{enumitem}
\usepackage{hyperref}

\geometry{margin=2.5cm}
\hypersetup{colorlinks=true,linkcolor=blue!60!black,urlcolor=blue!60!black}

\newcommand{\N}{\mathbb{N}}
\newcommand{\R}{\mathbb{R}}
\newcommand{\E}{\mathbb{E}}
\newcommand{\Prob}{\mathbb{P}}
\newcommand{\F}{\mathcal{F}}
\newcommand{\G}{\mathcal{G}}
\newcommand{\indep}{\perp\!\!\!\perp}
\newcommand{\one}{\mathbf{1}}

\theoremstyle{plain}
\newtheorem{thm}{Theorem}
\newtheorem{lem}[thm]{Lemma}
\newtheorem{cor}[thm]{Corollary}

\theoremstyle{definition}
\newtheorem{example}[thm]{Example}

\theoremstyle{remark}
\newtheorem*{remark}{Remark}

\title{The i.i.d.\ indicator admission \\
  \large Past-determined target meets independent uniform draw}
\author{ENS/Polytechnique 2026 Math A --- Lean 4 formalization\\
  Companion note (revised after external-LLM verification)}
\date{\today}

\begin{document}
\maketitle

\begin{abstract}
This note explains the probabilistic lemma admitted in the project as
\verb|iIndepFun_iIdentDistrib_uniformIndic_pastDep|. The earlier
draft of this note framed it as ``Williams \S 9.7'', citing
Williams, \emph{Probability with Martingales}, \S 9.7. \textbf{That
framing is misleading.} An external LLM, asked to prove the
underlying claim with all named-theorem invocations explicitly
forbidden, produced a direct $\sim$10-line counting argument
using only definitions of independence + uniformity + the
admitted local tree structure. No graduate student would write
``by Williams \S 9.7'' in their exam booklet. We give the direct
argument here, identify the specific Mathlib API gap that forces
the Lean admission anyway, and document its two project uses.
\end{abstract}

\section{Statement}

Throughout, $(\Omega, \F, \Prob)$ is a probability space and $S$ a
finite set of cardinality $|S| = n$. Equip $S$ with the discrete
$\sigma$-algebra. Let $(Y_k)_{k \geq 0}$ be a sequence of independent,
identically distributed random variables, each uniform on $S$:
\[
  \Prob(Y_k = s) = \frac{1}{n} \qquad \text{for every } s \in S.
\]
Define the natural filtration $\G_k := \sigma(Y_0, Y_1, \dots, Y_{k-1})$
for $k \geq 1$, with $\G_0 := \{\emptyset, \Omega\}$ trivial.

\begin{thm}[Williams \S 9.7, restated for our setting]
\label{thm:main}
Let $c \in \{0, 1, \dots, n\}$ and let $A_k : \Omega \to 2^S$ be a
sequence of $\G_k$-measurable random subsets of $S$ with constant
cardinality $|A_k(\omega)| = c$ for every $\omega \in \Omega$.
Define the indicator sequence
\[
  f_k(\omega) := \one_{\{Y_k(\omega) \in A_k(\omega)\}}, \qquad k \geq 0.
\]
Then $(f_k)_{k \geq 0}$ is an i.i.d.\ sequence of Bernoulli random
variables with parameter $c/n$.
\end{thm}

\noindent The two project applications correspond to specific choices of $S$,
$Y_k$, and $A_k$; see \S\ref{sec:applications}.

\section{The direct counting proof (what a grad student would write)}

The proof reduces to counting on the finite-product measure
$\Prob_M := \mathrm{Unif}(S)^{\otimes M}$ on $S^M$, where
$M = \max(I) + 1$ for the index set $I$ of interest.

\begin{proof}[Direct proof of Theorem~\ref{thm:main}]
Fix any finite $I = \{i_1 < \dots < i_r\} \subset \N_0$ and any
assignment $\epsilon : I \to \{0, 1\}$. We show
\[
  \Prob\bigl(\forall i \in I,\ f_i = \epsilon_i\bigr)
    = \prod_{i \in I} p_{\epsilon_i},
  \qquad p_1 := c/n,\ p_0 := 1 - c/n.
\]

The event depends only on the first $M$ coordinates, so it suffices
to count realizing tuples in $S^M$. For each prefix length
$\ell \in \{0, 1, \dots, M-1\}$, given a partial realization
$(y_0, \dots, y_{\ell-1}) \in S^\ell$ that satisfies the constraints
$f_i = \epsilon_i$ for all $i \in I \cap [0, \ell)$, count the
number $c_\ell$ of admissible choices for $y_\ell$:
\begin{itemize}
  \item If $\ell \notin I$: no constraint at step $\ell$; all $n$
    values of $y_\ell$ are admissible. So $c_\ell = n$.
  \item If $\ell \in I$ with $\epsilon_\ell = 1$: need $y_\ell \in
    A_\ell(y_0, \dots, y_{\ell-1})$, which has cardinality $c$. So
    $c_\ell = c$.
  \item If $\ell \in I$ with $\epsilon_\ell = 0$: need $y_\ell \notin
    A_\ell(\dots)$, cardinality $n - c$. So $c_\ell = n - c$.
\end{itemize}

Crucially, $c_\ell$ does \textbf{not} depend on the partial realization
$(y_0, \dots, y_{\ell-1})$; only on whether $\ell \in I$ and the
desired $\epsilon_\ell$. (The set $A_\ell(\dots)$ varies with the
prefix, but its \emph{size} is the constant $c$ by hypothesis.)

By induction on $\ell$ from $0$ to $M-1$, the total count is
$N = \prod_{\ell=0}^{M-1} c_\ell$. The probability of the event is
\[
  \frac{N}{n^M} = \prod_{\ell=0}^{M-1} \frac{c_\ell}{n}
    = \prod_{\ell \in I} \frac{c_{\ell}}{n}
    = \prod_{i \in I} p_{\epsilon_i}
\]
(steps $\ell \notin I$ contribute factor $1$). This establishes
mutual independence and identical Bernoulli$(c/n)$ distribution.
\end{proof}

\begin{remark}[Length and depth of the argument]
The substantive content is the bullet list and the
``$c_\ell$ does not depend on the prefix'' observation. That's two
sentences. Combined with the inductive product formula, the proof is
$\sim 10$ lines. There is no abstract conditional-expectation step.
The argument scales to any finite $I$ and any past-dependence
structure as long as the cardinality $|A_\ell| = c$ is constant.
\end{remark}

\section{An equivalent, more abstract framing (for context)}

The same fact is sometimes packaged as a corollary of a
measure-theoretic lemma:

\begin{lem}\label{lem:cond-law}
For every $k \geq 0$ and every $\G_k$-measurable event
$E \subseteq \Omega$,
\[
  \Prob(\{f_k = 1\} \cap E) = \frac{c}{n}\, \Prob(E).
\]
\end{lem}

\begin{proof}
The set $A_k$ is $\G_k$-measurable and takes values in the finite
collection $\binom{S}{c}$ of $c$-subsets of $S$. Decompose $E$ by
the value of $A_k$:
\[
  E = \bigsqcup_{T \in \binom{S}{c}} \big(E \cap \{A_k = T\}\big).
\]
Each piece $E_T := E \cap \{A_k = T\}$ is $\G_k$-measurable. On $E_T$,
the event $\{f_k = 1\} = \{Y_k \in T\}$. By independence of $Y_k$ from
$\G_k$ (which holds since $\G_k = \sigma(Y_0, \dots, Y_{k-1})$ and the
$Y_j$ are independent),
\[
  \Prob(\{Y_k \in T\} \cap E_T)
    = \Prob(Y_k \in T) \cdot \Prob(E_T)
    = \frac{|T|}{n} \cdot \Prob(E_T)
    = \frac{c}{n} \cdot \Prob(E_T).
\]
Sum over $T$:
\[
  \Prob(\{f_k = 1\} \cap E)
    = \sum_T \Prob(\{Y_k \in T\} \cap E_T)
    = \frac{c}{n} \sum_T \Prob(E_T)
    = \frac{c}{n} \cdot \Prob(E). \qedhere
\]
\end{proof}

\begin{cor}\label{cor:f-indep}
$f_k$ is independent of $\G_k$, and marginally $f_k \sim
\mathrm{Bernoulli}(c/n)$.
\end{cor}

\begin{proof}
Take $E = \Omega$ in Lemma~\ref{lem:cond-law}: $\Prob(f_k = 1) = c/n$
gives the marginal. For independence, we must show
$\Prob(\{f_k = a\} \cap E) = \Prob(f_k = a)\,\Prob(E)$ for $a \in
\{0, 1\}$ and any $\G_k$-measurable $E$. The case $a = 1$ is
Lemma~\ref{lem:cond-law}; the case $a = 0$ follows by complementation:
\[
  \Prob(\{f_k = 0\} \cap E) = \Prob(E) - \Prob(\{f_k = 1\} \cap E)
    = \Big(1 - \frac{c}{n}\Big)\,\Prob(E). \qedhere
\]
\end{proof}

\begin{proof}[Proof of Theorem~\ref{thm:main}]
We have established that each $f_k$ has marginal
$\mathrm{Bernoulli}(c/n)$, so the family is identically distributed.
For mutual independence, observe that $f_0, f_1, \dots, f_{k-1}$ are
$\G_k$-measurable (since $f_j$ is determined by $Y_0, \dots, Y_j$ and
$j \leq k - 1$ implies $\sigma(f_j) \subseteq \G_k$). By
Corollary~\ref{cor:f-indep}, $f_k$ is independent of $\G_k$, hence in
particular of $(f_0, \dots, f_{k-1})$. Inductively, the family
$(f_k)_{k \geq 0}$ is mutually independent. Combined with identical
distribution, the family is i.i.d.
\end{proof}

\begin{remark}
The hypothesis ``$|A_k|$ constant'' is essential. If $|A_k|$ were a
genuine random variable (say, $\G_k$-measurable but not constant),
then the conditional probability $\Prob(f_k = 1 \mid \G_k) = |A_k|/n$
would still be a deterministic function of the past, but \emph{not}
constant; $f_k$ would be conditionally Bernoulli with a random
parameter, and the marginal would involve averaging $|A_k|/n$ over
the past. Independence from the past would fail in general.
\end{remark}

\begin{lem}[``Constant conditional probability'' lemma]
\label{lem:cond-law}
For every $k \geq 0$ and every $\G_k$-measurable event
$E \subseteq \Omega$,
\[
  \Prob(\{f_k = 1\} \cap E) = \frac{c}{n}\, \Prob(E).
\]
\end{lem}

This is provable as before via decomposition over the finite values
of $A_k$. From here mutual independence follows by a tower-of-induction
argument, exactly as Williams \S 9.7 packages it.

\textbf{The framing question.} Both the direct counting proof and
the abstract lemma~\ref{lem:cond-law} are correct and equivalent.
The direct counting is what a graduate student would write in 10
lines on their exam booklet; the abstract lemma is what a textbook
records for general use. \emph{Neither is harder than the other; they
encode identical content at different abstraction levels.}

\section{Why ``the centred indicators are not i.i.d.\ '' is wrong}

A common heuristic, repeated in the project's earlier draft (errata
note E3 of \verb|RandomWalk.lean|), goes:

\begin{quote}\itshape
The centred indicator $\xi_i := \one_{\{Y_i \in A_i\}} - c/n$
\textbf{depends on the past} (through $A_i$), hence the family
$(\xi_i)$ is not independent.
\end{quote}

This is incorrect. ``Depends on the past'' is a statement about
\emph{measurability}, not about \emph{independence}. The
$\xi_i$ is $\sigma(Y_0, \dots, Y_i)$-measurable, true. But
Theorem~\ref{thm:main} says the family is nevertheless i.i.d.\
because the conditional law given the past is \emph{constant}:
``constant conditional law'' is a measure-theoretic property strictly
stronger than ``measurable function of past'' and strictly weaker
than ``deterministic''.

The constant-cardinality hypothesis $|A_k| = c$ is exactly the
content the heuristic misses.

\section{Project applications}\label{sec:applications}

In both applications, $S = \{a, b, a^{-1}, b^{-1}\}$ is the
generating set of $F_2$, equipped with the uniform measure
$Z_{\mathrm{unif}}$ ($n = 4$), and $Y_k$ is the $k$-th step of the
random walk on $F_2$.

\subsection{Q42 --- binomial identification}

The exam asks to show that
\[
  \frac{b_\varphi(X_n) + n}{2} \sim \mathcal{B}(n, 3/4).
\]
Write the LHS as $\sum_{i<n} \xi_i$ where
\[
  \xi_i := \frac{b_\varphi(X_{i+1}) - b_\varphi(X_i) + 1}{2}.
\]
By the admitted $1+3$ neighbour structure of the Busemann function,
exactly $1$ of the $4$ generators decreases $b_\varphi$ by~$1$ and $3$
increase it by~$1$. So $\xi_i = \one_{\{Y_i \in A_i\}}$ where
\[
  A_i(\omega) := \{ s \in S : b_\varphi(X_i(\omega) \cdot s)
    = b_\varphi(X_i(\omega)) + 1 \}
\]
has constant cardinality $|A_i| = 3$. By Theorem~\ref{thm:main},
$(\xi_i)$ is i.i.d.\ Bernoulli$(3/4)$, hence
$\sum_{i<n} \xi_i \sim \mathcal{B}(n, 3/4)$.

The exam's admitted Hoeffding inequality (preamble) then gives the
required tail bound, and Borel--Cantelli yields the SLLN
$b_\varphi(X_n)/n \to 1/2$ a.s.

\subsection{Q43 --- cancellation coupling}

The exam asks to show $|X_n|/n \to 1/2$ a.s. Writing $|X_n| = n - 2
C_n$ where $C_n$ counts cancellations, the proof reduces to $C_n/n
\to 1/4$ a.s. Define
\[
  \mathrm{coupled}_k := \one_{\{Y_k = (s_k)^{-1}\}}
\]
where $s_k$ is the last letter of $X_k$ (with arbitrary default at
identity). The set $A_k = \{(s_k)^{-1}\} \subseteq S$ has constant
cardinality $1$. By Theorem~\ref{thm:main}, $(\mathrm{coupled}_k)$ is
i.i.d.\ Bernoulli$(1/4)$. Mathlib's strong law gives $C_n/n \to 1/4$
a.s.\ (after a coupling argument that handles the random visits to
identity, where the formal $\mathrm{coupled}_k$ disagrees with the
true cancellation indicator on a finite set a.s.).

\subsection{Common structure}

Both applications fit Theorem~\ref{thm:main} with $|S| = 4$ and
constant $c \in \{1, 3\}$. They are the only places in the project
where the admission is invoked. Wave 23C consolidated them under a
single generic Lean axiom matching Theorem~\ref{thm:main}'s
statement.

\section{Necessity}

\subsection{Q42 cannot be done without it at graduate level}

The exam asks the \emph{student} to derive the binomial law of
$\sum \xi_i$. The derivation requires showing the family is i.i.d.
This is exactly the content of Theorem~\ref{thm:main}; there is no
shorter route.

Alternative: admit the binomial law directly. This replaces a single
general admission with a more specific one and does not help any
other question.

\subsection{Q43 alternative routes}

The cleanest grad-level proof is the i.i.d.\ coupling above. A
geometric ``$4$-ray'' argument exists but invokes either:
\begin{itemize}
  \item the same i.i.d.\ structure on the Busemann increments
    (no saving), or
  \item more elaborate facts about random walks on non-amenable groups
    (Kesten, Woess, Ch.~8), strictly stronger than what the exam
    expects.
\end{itemize}
The i.i.d.\ coupling using Theorem~\ref{thm:main} is the minimal
graduate-level argument. The errata document confirms this (Q43
\S(b), corrected route).

\section{The Lean obstacle (this is purely about Mathlib API, not the math)}

The direct counting proof (\S 2) and the abstract framing (\S 3) are
both pure mathematics and complete. Why, then, do we admit the
statement in Lean?

The answer is that Lean's Mathlib library has a specific gap:
\[
  \mu = \mathrm{Measure.infinitePi}(\mu_S)_{k \in \N}
\]
into ``finite prefix $\otimes$ infinite suffix'' over $\N$:
\[
  \mu = \mathrm{Measure.pi}\bigl(\mu_S^{[0,N)}\bigr)
    \otimes \mathrm{Measure.infinitePi}(\mu_S)_{k \geq N}.
\]
Mathlib has the finite-prefix split for $\mathrm{Measure.pi}$ over
\verb|Iic b| ranges (\verb|pi_prod_map_IicProdIoc|), but not for
$\mathrm{Measure.infinitePi}$ over $\N$. Filling this gap (\~200~LOC
of pure measure-theoretic plumbing, no probability or
graph-theoretic content) would dissolve the Lean admission and turn
Theorem~\ref{thm:main} into a Lean theorem.

This is a one-time Mathlib investment with multiple downstream
benefits beyond this project.

Once the prefix-suffix split is in place, the Lean proof of
Theorem~\ref{thm:main} is the same direct counting argument as
\S 2: project to $S^M$, count $N = \prod c_\ell$, divide by $n^M$.
$\sim 200$ LOC of measure-theoretic plumbing for a 10-line
mathematical argument.

\textbf{The honest situation:} the Lean admission is misleadingly
named after Williams \S 9.7 because that's a convenient textbook
citation. The mathematical content is undergraduate-tier counting;
the Lean blocker is purely the absence of a finite-prefix split for
\verb|Measure.infinitePi|. Future readers should understand: this is
not a deep probability admission, it is a Lean infrastructure gap.

\section{References}

\begin{itemize}[leftmargin=*]
  \item D.~Williams, \emph{Probability with Martingales}, Cambridge
    University Press, 1991. \S 9.7 (``Conditional expectations:
    further properties'').
  \item A.~Klenke, \emph{Probability Theory: A Comprehensive Course},
    Springer, 2014, 2nd~ed. \S 5.3 (``Conditional probabilities and
    conditional expectations'').
  \item D.~Bertsekas, J.~Tsitsiklis, \emph{Introduction to
    Probability}, Athena Scientific, 2008. \S 1.5--1.6 covers the
    finite-state version (``constant conditional probability $\Rightarrow$
    independence'') as a worked exercise.
\end{itemize}

\end{document}
